\documentclass[aps,12pt,floatfix,longbibliography,a4paper]{revtex4-2}
%twocolumn, 

\usepackage{graphicx}
\usepackage{epstopdf}
\usepackage{amssymb}
\usepackage{hyperref}
\usepackage[normalem]{ulem}
\usepackage{color}
\usepackage{amsmath,bm}
\usepackage[utf8]{inputenc}
\usepackage{xcolor}
\usepackage{physics}
\usepackage{mathcomp}
\usepackage{multirow}
\usepackage{comment}
\usepackage{makecell}

\usepackage[labeled]{multibib}
%\usepackage[resetlabels,labeled]{multibib}
\usepackage{placeins}

%\usepackage{cmbright}
%\DeclareFontShape{OT1}{cmss}{m}{it}{<->ssub*cmss/m/sl}{}
%\renewcommand{\rmdefault}{cmss}
%\renewcommand{\sfdefault}{cmss}
%\setlength{\parindent}{0em}
%\setlength{\parskip}{1em}

\newcommand{\sgn}{\mathop{\mathrm{sgn}}\nolimits}
\newcommand{\LNO}{La$_3$Ni$_2$O$_7$}
\newcommand{\Tc}{T$_\mathrm{c}$}


\newcommand{\RO}[1]{\textcolor{blue}{#1}}

\newcounter{hoge}
\makeatletter
\@addtoreset{figure}{hoge}%%%(*2)
\makeatother
\makeatletter
\@addtoreset{equation}{hoge}%%%(*2)
\makeatother
\makeatletter
\@addtoreset{section}{hoge}%%%(*2)
\makeatother
\makeatletter
\@addtoreset{table}{hoge}%%%(*2)
\makeatother

\begin{document}

\title[]{Polar, checkerboard charge order in bilayer nickelate \LNO{}}

%%=============================================================%%

\author{Ryo Misawa$^{1}$}
\email{misawann6@g.ecc.u-tokyo.ac.jp}
\author{Shunsuke Kitou$^{2}$}
\author{Jian-Ping Sun$^{3,4}$}
\author{Yingpeng Yu$^{3,4}$}
\author{Chihaya Koyama$^{2}$}
\author{Yuiga Nakamura$^{5}$}
\author{Taka-hisa Arima$^{2,6}$}
\author{Jin-Guang Cheng$^{3,4}$}
\author{Max Hirschberger$^{1}$}
\email{hirschberger@ap.t.u-tokyo.ac.jp}

\affiliation{$^{1}$Department of Applied Physics, The University of Tokyo, Bunkyo, Tokyo 113-8656, Japan}
\affiliation{$^{2}$Department of Advanced Materials Science, The University of Tokyo, Kashiwa, Chiba 277-8561, Japan}
\affiliation{$^{3}$Beijing National Laboratory for Condensed Matter Physics, and Institute of Physics, Chinese Academy of Sciences, Beijing 100190, China}
\affiliation{$^{4}$School of Physical Sciences, University of Chinese Academy of Sciences, Beijing 100190, China}
\affiliation{$^{5}$Japan Synchrotron Radiation Research Institute (JASRI), SPring-8, Hyogo 679-5198, Japan}
\affiliation{$^{6}$RIKEN Center for Emergent Matter Science (CEMS), Wako, Saitama 351-0198, Japan}

\maketitle
\newpage

\begin{center}
\Large{Abstract}
\end{center}

Competing charge and spin orders are central to uncovering the nature of unconventional superconductivity. Here we utilize synchrotron X-ray diffraction on a high-quality single crystal to reveal the charge order of \LNO{} at ambient pressure, which competes with the high-temperature superconducting phase under pressure. Enabled by the high synchrotron photon flux and a large dynamic range, we resolve faint reflections---nearly four orders of magnitude weaker than the main Bragg reflections---that were overlooked in prior diffraction studies. This observation evidences a broken glide-mirror symmetry, leading to a polar crystal structure, rather than the widely used centrosymmetric structure model. The polarity is induced by checkerboard charge order on nickel sites in combination with octahedral tilting, reminiscent of bilayer manganese oxides. Our results provide a foundation for understanding phase competition and the mechanism of pressure-induced superconductivity in bilayer nickelates.

\newpage

\begin{center}
\Large{Main Text}
\end{center}
%%%%%%%%%%%%%%%%%%%%%%%%%%%%%
%%% SECTION
%%%%%%%%%%%%%%%%%%%%%%%%%%%%%
In the search for high-temperature superconductivity beyond the paradigm of Bardeen-Cooper-Schrieffer theory, transition-metal oxides based on copper, as well as iron pnictides and chalcogenides, have served as central material platforms~\cite{Keimer2015-zv,Stewart2011-qs,Fernandes2022-gu}. Typically, such unconventional superconductivity emerges alongside a variety of charge- and spin-ordered phases; these phases can compete or coexist with superconductivity~\cite{Keimer2015-zv,Fradkin2015-px}. More recently, high-temperature superconductivity was discovered in thin films of NdNiO$_2$ at ambient pressure~\cite{Li2019-jv}. Subsequently, pressure-induced superconductivity was observed in single crystals of the bilayer and trilayer Ruddlesden-Popper nickelate La$_{n+1}$Ni$_{n}$O$_{3n+1}$ ($n=2,3$)~\cite{Sun2023-sf,Zhang2024-uu,Zhu2024-rn}. Focusing on the bilayer nickelate \LNO{}, a majority of crystals are reported to crystallize in the orthorhombic $Amam$ structure at ambient pressure (Fig.~\ref{main:Fig1}\textbf{a}). Distinct structures built from monolayer and trilayer blocks are known polymorphs of \LNO{} (see Methods). The widely used $Amam$ structure is inferred largely from powder X-ray diffraction (XRD)~\cite{Zhang1994-vm,Taniguchi1995-lg,Ling2000-uz,Voronin2001-oc,Wang2024-qi} and from in-house single-crystal XRD~\cite{Wang2025-tg,Chen2024-dr,Wang2024-ii}. This structure model contains a single nickel site with the nominal valence of $2.5+$, corresponding to a mixed-valence state. Furthermore, a seminal powder XRD study reports a symmetry ascending to the tetragonal $I4/mmm$ structure under pressure~\cite{Wang2024-qi}. Around the critical pressure where the tetragonal structure is first observed at room temperature, La$_3$Ni$_2$O$_7$ shows superconductivity up to $80\,$K~\cite{Sun2023-sf,Zhang2024-uu}. These observations suggest that the ambient-pressure structure competes with the pressure-induced superconductivity, underlining the need to clarify the structural properties at ambient pressure with high experimental resolution. Here, we demonstrate polar charge order in \LNO{} at ambient conditions (Fig.~\ref{main:Fig1}\textbf{b}-\textbf{e}), directly evidenced by a broken glide-mirror symmetry in XRD experiments.

%%%%%%%%%%%%%%%%%%%%%%%%%%%%%
%%% SECTION
%%%%%%%%%%%%%%%%%%%%%%%%%%%%%
\bigskip
\textbf{Breaking of glide-mirror symmetry}\\
To precisely investigate the crystal structure of \LNO{} at ambient pressure, we perform synchrotron X-ray diffraction on a high-quality single crystal (Methods). We use an area detector with a dynamic range of $10^6$, enabling simultaneous detection of weak oxygen-derived intensities and much stronger reflections dominated by lanthanum and nickel. From systematic extinction of reflections in the observed diffraction patterns, we identify an orthorhombic $A$-base-centered lattice at and below room temperature, consistent with previous reports~\cite{Zhang1994-vm,Taniguchi1995-lg,Ling2000-uz,Voronin2001-oc,Wang2024-qi,Wang2025-tg}. However, there is evidence of symmetry breaking in the plane of reciprocal space described by the Miller indices $h0l$, which was not previously reported. Figure~\ref{main:Fig1}\textbf{f} shows the XRD diffraction pattern on this plane, taken at temperature $T=50\,$K. We resolve the systematic presence of $h0l$ ($h=$~odd) reflections. This corresponds to the absence of the $a$-glide plane, in clear disagreement with the reported space group $Amam$~\cite{Zhang1994-vm,Taniguchi1995-lg,Ling2000-uz,Voronin2001-oc,Wang2024-qi,Wang2025-tg}. Since these reflections are observed at and below room temperature (see Supplementary Note), we hereafter focus on the data at $T=50\,$K, at which the intensities are the strongest among the measured temperatures.

Quantitative evidence against the $Amam$ structure is provided by high-precision structure refinement, enabled by an excellent coverage of reciprocal space: $98\,\%$ of all available reflections, up to a resolution of of $0.3249\,$\AA, were measured by rotating the sample in fine steps, yielding the redundancy $16.296$ (Methods). Figure~\ref{main:Fig2}\textbf{a} compares the observed and calculated XRD intensities for the best refinement based on the $Amam$ structure, in terms of  $|F_\mathrm{obs}|^2$ and $|F_\mathrm{calc}|^2$. Here, $F_\mathrm{obs}$ and $F_\mathrm{calc}$ represent the observed and calculated structure factors, respectively. Even when ignoring the reflections that break the $a$-glide plane as symmetry-extinct in the refinement, the result shows a poorer agreement between $|F_\mathrm{obs}|^2$ and $|F_\mathrm{calc}|^2$ with reliability factor $R = 4.22\,\%$ and goodness-of-fit (GOF) of $7.13$ based on $5,542$ unique reflections; see Table~\ref{Tab:refinement} for other figures of merit of this refinement. The thermal ellipsoids of oxygen atoms show an anomalously flat shape in the refinement using $Amam$, further suggesting the need to lower the structural symmetry~(see Supplementary Note).

To identify the correct space group, we consider subgroups of the parent tetragonal structure $I4/mmm$; the reflection conditions uniquely determine the highest-symmetry space group as polar $Am2m$. Indeed, our refinement of XRD intensities using this structure yields much better agreement, with $R=1.36\,\%$ and $\text{GOF}=1.15$ based on $10,834$ unique reflections (Fig.~\ref{main:Fig2}\textbf{b}). The absence of the inversion center is further evidenced by the Flack parameter, i.e., the inequivalent occupation of inversion-paired domains in our single crystal: $0.63(2)$ and $0.37(2)$. Moreover, the refinement of oxygen occupancies shows no deviation from the stoichiometric value, and thermal ellipsoids at the oxygen sites are small and regular in shape (see Supplementary Note). These results indicate that the faint streaks in the XRD pattern (Fig.~\ref{main:Fig1}\textbf{f}), much weaker than the Bragg reflections, play only a minor role here, in contrast to the severe diffuse streaks arising from stacking disorder in a polymorph of \LNO{} (see Methods). Together, our analysis excludes the widely assumed $Amam$ structure and determines the proper structure as $Am2m$.

%%%%%%%%%%%%%%%%%%%%%%%%%%%%%
%%% SECTION
%%%%%%%%%%%%%%%%%%%%%%%%%%%%%
\bigskip
\textbf{Discovery of polar, checkerboard charge order in \LNO{}}\\
We now examine the refined structure in detail. Unlike the $I4/mmm$ and $Amam$ structure models, the proposed $Am2m$ structure contains two inequivalent nickel sites, as shown in Fig~\ref{main:Fig1}\textbf{b},\textbf{c}. These two sites have distinct Ni-O bond lengths, whose difference is rather large, $\sim 0.1~\text{\AA}$. We note that a recent report lowered symmetry to monoclinic $P2_1/m$ based on in-house single-crystal XRD, but found Ni–O bond-length differences of only $0.01$~\AA\ between inequivalent Ni sites~\cite{Li2026-go}. This is an order of magnitude smaller than in our synchrotron XRD refinement (Fig.~\ref{main:Fig1}\textbf{c}); it may be difficult to resolve weak intensities related to oxygen displacements in in-house XRD, without the high synchrotron photon flux and large dynamic range employed in our experiment. 

The large difference in the Ni-O bond lengths is readily explained by charge order on the nickel sites, a reasonable scenario for an oxide with a nominal nickel valence of $2.5+$. The valence on each nickel site can be estimated by the well-established bond-valence sum approach (Methods): the nickel site with a contracted octahedron has a valence of $+2.785(2)$, whereas the other nickel site has a valence of $+2.347(1)$. The absence of the $a$-glide symmetry is depicted in Fig.~\ref{main:Fig1}\textbf{c}, where the glide operation connects the two nickel sites with different valences.

The two nickel sites alternate in each layer and form a checkerboard pattern of charge. This checkerboard charge order, combined with octahedral tilting of the oxygen ligand environment, imparts the polarity to \LNO{}. Figure~\ref{main:Fig1}\textbf{d} illustrates a single polar domain as seen from the side; its inversion pair is obtained by swapping the charge-rich and charge-poor sites (Fig.~\ref{main:Fig1}\textbf{e}). Note that, without octahedral tilting, the checkerboard charge distribution within a single layer would result in full cancellation of the net polarity. Only when octahedral tilting introduces bond alternation along the same direction ($b$ axis) as the charge alternation, cancellation no longer occurs, and a macroscopic polarity arises.

%%%%%%%%%%%%%%%%%%%%%%%%%%%%%
%%% SECTION
%%%%%%%%%%%%%%%%%%%%%%%%%%%%%
\bigskip
\textbf{Discussion}\\
The checkerboard charge order discovered here is a plausible yet previously unrecognized state in the bilayer nickel oxide \LNO{} with mixed valence of nickel. Indeed, our charge order harks back to the structure of the bilayer manganese oxide Pr(Sr$_{1-x}$Ca$_x$)$_2$Mn$_2$O$_7$, which likewise exhibits electric polarization due to the alternation of charges and bonds~\cite{Tokunaga2006-bz}. Related, yet different, structures have been proposed in other bilayer Ruddlesden-Popper transition-metal oxides: out-of-plane polarity of Sr$_3$Co$_2$O$_7$ is theoretically linked to charge order~\cite{Zhou2026-uk,Huang2026-ax}, in-plane polarity in Ca$_3$Ru$_2$O$_7$ arises from rotation and tilting of octahedra~\cite{Lei2018-lg}, and hidden order in Sr$_3$Fe$_2$O$_7$ is ascribed to centrosymmetric checkerboard charge order by resonant X-ray diffraction~\cite{Kim2021-vw}. 
Since the ambient-pressure structure of \LNO{} turns out to be non-centrosymmetric, the pressure-induced transition to the $I4/mmm$ structure warrants re-examination using high-resolution experimental probes, beyond powder diffraction. By identifying the checkerboard charge order with broken inversion symmetry as a competing phase of pressure-induced superconductivity, this work reshapes the current understanding of the Ruddlesden-Popper nickel oxides and calls for further experimental and theoretical studies of the competition between charge order and superconductivity.

\newpage
\bibliography{sn-bibliography}

\newpage

%%===========================================================================================%%
%% MAIN TEXT FIGURES
%%===========================================================================================%%
\clearpage

\begin{center}
\Large{Main Text Figures}
\end{center}
\FloatBarrier
\vspace{5mm}
\clearpage
\begin{figure}[ht]%
\centering
\includegraphics[width=1.0\textwidth]{Fig1.pdf}
\caption{\textbf{Polar charge order in bilayer nickelate \LNO{} at ambient pressure.}
\textbf{a}, Centrosymmetric, orthorhombic $Amam$ structure model of \LNO{}, previously reported in the literature. There is only a single independent nickel site, surrounded by gray oxygen octahedra. \textbf{b}, Polar orthorhombic $Am2m$ structure proposed in the present work. This model contains two inequivalent nickel sites with smaller (orange) and larger (blue) NiO$_6$ octahedral environments. The orange and blue lines indicate shorter and longer Ni-O bonds. \textbf{c}, Broken $a$-glide symmetry in the $Am2m$ structure. The glide operation maps a nickel site to the other independent site; the symmetry is broken if the sites are inequivalent. Broken glide-mirror plane indicated in gray. \textbf{d}, Polarity $P$ induced by the charge alternation along the $b$ axis in combination with the bond alternation due to octahedral tilting around the $a$ axis. The nickel sites with orange and blue octahedra carry a valence of $+2.785(1)$ and $+2.347(1)$, respectively, as estimated from the valence-bond sum (Methods). The tilting angles are $\theta_1 = 192.183(16)^\circ$ and $\theta_2=170.945(14)^\circ$. \textbf{e}, Another polar domain obtained by swapping charges in panel~\textbf{d}. \textbf{f}, X-ray diffraction pattern on the $h0l$ plane at $50~$K. Here, $h0l$ ($h=-1$) reflections associated with the broken $a$-glide plane are highlighted by the black rectangle.
}
\label{main:Fig1}
\end{figure}

\clearpage
\begin{figure}[ht]%
\centering
\includegraphics[width=0.5\textwidth]{fig2.pdf}
\caption{\textbf{Structure refinement based on synchrotron X-ray diffraction (XRD).}
\textbf{a}, Refinement of XRD data at $T = 50\,$K assuming the $Amam$ structure model. Observed and calculated intensities, expressed in terms of $|F_\mathrm{obs}|^2$ and $|F_\mathrm{calc}|^2$, show a poor agreement in the regime of weak intensity. Here, $F_\mathrm{obs}$ and $F_\mathrm{calc}$ denote the observed and calculated structure factors, respectively. The poor agreement of model and experiment is quantified by the reliability factor $R = 4.22\,\%$ and goodness-of-fit (GOF) of $7.13 \gg 1$. \textbf{b}, Refinement assuming the proposed $Am2m$ structure. $|F_\mathrm{obs}|^2$ and $|F_\mathrm{calc}|^2$ show a much better agreement with $R = 1.28\,\%$ and $\text{GOF}=1.07$. Inset focuses on $h0l$ ($h=\text{odd}$) reflections (orange circles) associated with the broken $a$-glide.
}
\label{main:Fig2}
\end{figure}

\clearpage
\newpage
\begin{table}[ht]
\centering
\caption{\textbf{Refinement statistics for \textit{Amam} and \textit{Am}2\textit{m} models of \LNO{}.} Oxygen occupancies are fixed to $1$, since the refinement results in the stoichiometric values for the $Am2m$ structure; see Supplementary Note.}
\label{Tab:refinement}
\setlength{\tabcolsep}{6pt}
\begin{tabular}{c c c}
\hline
Space group & $Amam$ (No.~63) & $Am2m$ (No.~38)\\
\hline
Temperature (K) & \multicolumn{2}{c}{$50$} \\
Wavelength ($\text{\AA}$) & \multicolumn{2}{c}{$0.343008$} \\
$a$ ($\text{\AA}$) & \multicolumn{2}{c}{$5.37640(10)$} \\
$b$ ($\text{\AA}$) & \multicolumn{2}{c}{$5.45650(10)$} \\
$c$ ($\text{\AA}$) & \multicolumn{2}{c}{$20.5013(4)$} \\
$\alpha = \beta = \gamma$ ($^\circ$) & \multicolumn{2}{c}{$90$} \\
Cell volume ($\text{\AA}^{3}$) & \multicolumn{2}{c}{$601.43(2)$} \\
$Z$ & \multicolumn{2}{c}{$4$} \\
$F(000)$ & \multicolumn{2}{c}{$1132$} \\
$d_\text{min}$ (\AA) & \multicolumn{2}{c}{$0.3249$} \\
Completeness (\%) & $98$ & $97$ \\
Redundancy & $16.296$ & $8.527$ \\
$N_\mathrm{total}$ & \multicolumn{2}{c}{$92383$} \\
$N_\mathrm{unique}$ & $5542$ & $10834$ \\
\hline
Inversion twin fractions & Not applicable & $0.63(2)$, $0.37(2)$ \\
$N_\mathrm{parameters}$ & $37$ & $72$ \\
$R_1\ (I>3\sigma(I)\ /\ \text{all})$ & $4.19\ \%$ / $4.22\ \%$ & $1.24\ \%$ / $1.28\ \%$ \\
$wR_2\ (I>3\sigma(I)\ /\ \text{all})$ & $10.39\ \%$ / $10.40\ \%$ & $1.73\ \%$ / $1.75\ \%$ \\
GOF $(I>3\sigma(I)\ /\ \text{all})$ & $7.20$ / $7.13$ & $1.07$ / $1.07$ \\
\hline
\end{tabular}
\end{table}

\clearpage
\newpage

%%%%%%%%%%%%%%%%%%%%%%%%%%%%%%%%%%%%%%%%%%%%%%%%%%%%%%%%%%%%%%%%%%%%%%%
%%%%%%%%%%%%%%%%%%%%%%%%%%%%%%%%%%%%%%%%%%%%%%%%%%%%%%%%%%%%%%%%%%%%%%%
%%				METHODS
%%%%%%%%%%%%%%%%%%%%%%%%%%%%%%%%%%%%%%%%%%%%%%%%%%%%%%%%%%%%%%%%%%%%%%%
%%%%%%%%%%%%%%%%%%%%%%%%%%%%%%%%%%%%%%%%%%%%%%%%%%%%%%%%%%%%%%%%%%%%%%%
\begin{center}
\Large{Methods}
\end{center}

\textbf{Single crystal growth}\\
Single crystals of \LNO{} were grown by a molten salt flux evaporation method at ambient pressure~\cite{Li2026-go}. High-purity La$_2$O$_3$ ($99.99~\%$) and NiO were used as starting materials. La$_2$O$_3$ was dried at $1000~^\circ$C overnight, and the oxide powders were weighed and ground. The anhydrous K$_2$CO$_3$ was used as flux and then mixed with the oxide powders in a flux-to-solute mass ratio of $15:1$. To prevent moisture absorption, all the processes were performed inside a glovebox. The mixture was placed in an alumina crucible that was covered with a lid to control the evaporation rate. The crystals were grown in the furnace at a temperature of $1000\text{-}1050~^\circ$C for $72$ hours while evaporating the flux gradually, and then were cooled to room temperature naturally. The crystals were extracted by soaking the mixture in deionized water. 
\bigskip

\textbf{Synchrotron X-ray diffraction}\\
Synchrotron XRD measurements were performed at beamline BL$02$B$1$ of the SPring-$8$ synchrotron radiation facility (Japan). A single crystal of \LNO{} with a size of $\sim 50\,\mu$m was employed in the measurements. A N$_2$-gas-blowing device was used for temperatures between $100\,$K and $300\,$K, and a He-gas-blowing device for temperatures below $100\,$K. Bragg reflections were recorded using a CdTe PILATUS area detector and the \textsc{CrysAlisPro} software~\cite{Agilent-Technologies-Ltd2014-hl} with the fine-slice acquisition method. In this method, the reciprocal space was scanned in steps of $\Delta \omega = 0.01^\circ$. Symmetry-equivalent reflections were averaged, and structure refinements were performed with \textsc{Jana2006}~\cite{Petricek2014-pc}.
\bigskip

\textbf{Comparison of structure models in \LNO{}}\\
The highest-symmetry structure of the well-known bilayer Ruddlesden-Popper transition metal oxides corresponds to a monotonic stacking of bilayers (sequence: $2222$) with tetragonal space group $I4/mmm$. As a subgroup of $I4/mmm$, most studies of \LNO{} focused on the $Amam$ structure at ambient pressure~\cite{Zhang1994-vm,Taniguchi1995-lg,Ling2000-uz,Voronin2001-oc,Wang2024-qi,Wang2025-tg}, which is corrected to $Am2m$ in the present work. As discussed in the Main Text, $Am2m$ allows for two independent nickel sites due to charge ordering -- a major difference to $Amam$. On the other hand, polymorphism is also reported in some seminal works on \LNO{} crystals, depending on crystal growth conditions, for example, when using the floating-zone method. The polymorphs with space groups of $Cmmm$~\cite{Chen2024-dr,Wang2024-ii} and $Fmmm$~\cite{Puphal2024-pw} are characterized by a stacking of monolayers and trilayers ($1313$), instead of $2222$. The $Fmmm$ polymorph is known to exhibit stacking disorder, as evidenced by pronounced streaks in synchrotron XRD measurements~\cite{Puphal2024-pw}. Another stacking pattern of $1212$ is also reported~\cite{Li2024-vg}. 
In the present study, we focus on a high-quality single crystal with uniform bilayer stacking sequence, consistent with the best crystals reported in the literature~\cite{Zhang1994-vm,Taniguchi1995-lg,Ling2000-uz,Voronin2001-oc,Wang2024-qi,Wang2025-tg}. The suppressed thermal ellipsoids and the faintness of streaks in our XRD pattern (Fig.~\ref{main:Fig1}\textbf{f}) indicate that this bilayer structure is robust against stacking disorder, unlike previous reports on polymorphs of \LNO{}.

\newpage

\clearpage
\textbf{Acknowledgments}\\
We thank Max~T.~Birch and M.~Kriener for discussions. This work was supported by JSPS KAKENHI Grants No. JP23H05431, No. JP24H01607, JP24K17006, and 24H01644, as well as JST CREST Grant Nos. JPMJCR1874 and JPMJCR20T1 (Japan) and JST FOREST Grant No. JPMJFR2238, and No. JPMJFR2362 (Japan). It was also supported by JST as part of Adopting Sustainable Partnerships for Innovative Research Ecosystem (ASPIRE), Grant Number JPMJAP2426. M.H. is supported by the Deutsche Forschungsgemeinschaft (DFG, German Research Foundation) via Transregio TRR 360 – 492547816. J.P.S. and J.G.C. are supported by the National Key R\&D Program of China (2023YFA1406100) and the National Natural Science Foundation of China (12025408 and 12522407). The synchrotron single-crystal X-ray experiments were performed at BL02B1 of SPring-8 with the approval of RIKEN (Proposal Nos. 2025B1928, 2025B1855, and 2025B2178). 

\bigskip
\textbf{Author contributions}\\
M.H. conceived the project. J.P.S., Y.Y., and J.G.C. grew and characterized the single crystals. R.M., C.K., S.K., and Y.N. performed synchrotron XRD measurements. R.M. analyzed the XRD data in consultation with S.K and T.-h.A. R.M. wrote the manuscript in collaboration with M.H., S.K., and T.-h.A.; all authors discussed the results and commented on the manuscript.\\

\textbf{Competing interests}\\
The authors declare no competing interests.\\

\end{document}